\section{The integral monodromy of the period system}\label{sec:monodromy}

The integral period system is the part of the construction that can be determined completely by algebra, and it is used elsewhere: the author's higher-torus update attaches a character to its peripheral monodromy (Section~\ref{sec:later}), and the Erdős--Straus workbench builds covers of the projective line from its dual action (Section~\ref{sec:overlaps}). This section determines the monodromy group exactly and explains how these two uses fit it. All results of this section are new, and each is \textsf{proved here}; they were checked by computer and re-derived by an independent referee (Appendix~\ref{app:verification}). The existence question for the complex structure on $S^6$ itself is not addressed.

\subsection{Setting}

The period system is taken from the manuscript \cite{AlpogeS6}, as reconstructed in the record's file \fn{03\_}, §2. Let $V=\ZZ^4$ with basis $(\gamma,u,w,\delta)$. The local monodromies at the two elliptic points are, by formula (2.2) of the manuscript,
\begin{align*}
T_1:&\quad \gamma\mapsto\gamma,\quad u\mapsto-u-w,\quad w\mapsto-6\gamma+u,\quad \delta\mapsto2\gamma+u+w+\delta,\\
T_2:&\quad \gamma\mapsto\gamma,\quad u\mapsto6\gamma+w,\quad w\mapsto-u,\quad \delta\mapsto-3\gamma+u+\delta,
\end{align*}
and the monodromy at the cusp is $T_0=(T_1T_2)^{-1}$. By the manuscript's Lemma~2.2, $T_1^3=T_2^4=I$, and $T_0=I+N$ with $N^2=0$, $N\gamma=Nu=0$, $Nw=-u$ and $N\delta=\gamma$, so $\ker N=\operatorname{im}N=\langle\gamma,u\rangle$.

By the manuscript's Lemma~2.8 the group of $\langle T_1,T_2\rangle$-invariant alternating forms is infinite cyclic, generated by $Q_0$ with $Q_0(\gamma,\delta)=1$ and $Q_0(u,w)=6$, all other pairings of basis vectors being zero. Its elementary divisors are $(1,1,6,6)$, so $Q_0$ has type $(1,6)$. We write $G=\langle T_1,T_2\rangle$ for the monodromy group and
\[
G_{\ZZ}=\{g\in Sp(Q_0,\ZZ):\ g\gamma=\gamma\}
\]
for the stabiliser of $\gamma$. Every $T_i$ fixes $\gamma$ and preserves $Q_0$, so $G\subseteq G_{\ZZ}$.

\subsection{The stabiliser of \texorpdfstring{$\gamma$}{gamma}}

For $M\in SL_2(\ZZ)$ let $L_M$ act by $M$ on the plane $\langle u,w\rangle$, in the basis $(u,w)$, and fix $\gamma$ and $\delta$. For $p,r,s\in\ZZ$ let $h(p,r,s)$ be the map
\[
\gamma\mapsto\gamma,\qquad u\mapsto u-6r\gamma,\qquad w\mapsto w+6p\gamma,\qquad \delta\mapsto\delta+pu+rw+s\gamma.
\]

\begin{lemma}\label{lem:stab}
Every element of $G_{\ZZ}$ can be written uniquely as $L_Mh(p,r,s)$. The elements $h(p,r,s)$ form a group $\mathrm{Heis}(\ZZ)$ with the law
\[
h(p_1,r_1,s_1)\,h(p_2,r_2,s_2)=h\bigl(p_1+p_2,\ r_1+r_2,\ s_1+s_2+6(p_1r_2-p_2r_1)\bigr),
\]
and $L_Mh(p,r,s)L_M^{-1}=h(M(p,r),s)$. Hence $G_{\ZZ}=SL_2(\ZZ)\ltimes\mathrm{Heis}(\ZZ)$.
\end{lemma}
\begin{proof}
Let $g\in G_{\ZZ}$. Since $g$ fixes $\gamma$, it preserves $\gamma^\perp=\langle\gamma,u,w\rangle$. The form induced on $\gamma^\perp/\langle\gamma\rangle$ is $6\,(u^*\wedge w^*)$, so the induced map $M$ has determinant $1$. On $V/\gamma^\perp$, $g$ acts trivially, since $Q_0(\gamma,gx)=Q_0(g^{-1}\gamma,x)=Q_0(\gamma,x)$. Then $L_M^{-1}g$ acts trivially on $\langle\gamma\rangle$, on $\gamma^\perp/\langle\gamma\rangle$ and on $V/\gamma^\perp$. Such a unipotent map has the form $u\mapsto u+a_u\gamma$, $w\mapsto w+a_w\gamma$, $\delta\mapsto\delta+pu+rw+s\gamma$, and it preserves $Q_0$ exactly when $a_u=-6r$ and $a_w=6p$. The law and the conjugation formula are matrix products.
\end{proof}

Two consequences are used repeatedly. Inside $\mathrm{Heis}(\ZZ)$, conjugation acts by
\[
h(p_0,r_0,s_0)\,h(p,r,s)\,h(p_0,r_0,s_0)^{-1}=h\bigl(p,r,s+12(p_0r-pr_0)\bigr).
\] The centre of $G_{\ZZ}$ is $\{h(0,0,s)\}=\langle t_\gamma\rangle$, where $t_\gamma=h(0,0,1)$ is the transvection $\delta\mapsto\delta+\gamma$, that is, $x\mapsto x+Q_0(\gamma,x)\gamma$. The subgroup $H(\ZZ)=\{h(p,r,6\kappa)\}$ carries the standard Heisenberg law $\kappa_1+\kappa_2+(p_1r_2-p_2r_1)$ and has index $6$ in $\mathrm{Heis}(\ZZ)$; so $G_{\ZZ}$ contains the Jacobi group $SL_2(\ZZ)\ltimes H(\ZZ)$ with index $6$, and is generated by it together with $t_\gamma$.

\subsection{The monodromy group is the kernel of a character of order 12}

Let $\operatorname{ab}:SL_2(\ZZ)\to\ZZ/12$ be the abelianisation, normalised by $\operatorname{ab}(T)=1$ for $T=\begin{psmallmatrix}1&1\\0&1\end{psmallmatrix}$. Then $\operatorname{ab}(S)=9$ for $S=\begin{psmallmatrix}0&-1\\1&0\end{psmallmatrix}$, $\operatorname{ab}(R)=4$ for $R=\begin{psmallmatrix}-1&1\\-1&0\end{psmallmatrix}$, and $\operatorname{ab}(-I)=6$. These values are forced: $SL_2(\ZZ)^{\mathrm{ab}}\cong\ZZ/12$ is generated by the class of $T$ (the record's file \fn{15\_} proves this by a Smith normal form in its Theorem~6.1); $S$ and $R$ have orders $4$ and $3$; and $\operatorname{ab}(S)+\operatorname{ab}(R)=\operatorname{ab}(SR)=\operatorname{ab}(T)$, because $SR=\begin{psmallmatrix}1&0\\-1&1\end{psmallmatrix}$ is conjugate to $T$. On $\mathrm{Heis}(\ZZ)$ put
\[
\varphi(h(p,r,s))=s-6q(p,r)\bmod12,\qquad q(p,r)=p^2+pr+r^2\bmod2,
\]
and on $G_{\ZZ}$ put $\Phi(L_Mh)=\operatorname{ab}(M)+\varphi(h)\bmod12$.

\begin{theorem}[the monodromy group]\label{thm:mono}
\begin{enumerate}[label=(\alph*)]
\item $\Phi:G_{\ZZ}\to\ZZ/12$ is a surjective homomorphism, and $G=\ker\Phi$.
\item The local monodromies decompose as $T_1=L_Rh(-1,0,2)$, $T_2=L_Sh(0,-1,-3)$ and $T_0=L_{T^{-1}}t_\gamma$. In particular the Levi image of $G$ is all of $SL_2(\ZZ)$.
\item $G$ is normal in $G_{\ZZ}$, and $G_{\ZZ}/G\cong\ZZ/12$ is cyclic, generated by the class of the central transvection $t_\gamma$. Moreover $G_{\ZZ}=G\cdot\langle t_\gamma\rangle$ and $G\cap\langle t_\gamma\rangle=\langle t_\gamma^{12}\rangle$.
\item $G\cap\mathrm{Heis}(\ZZ)=H':=\{h(p,r,s):s\equiv6(p^2+pr+r^2)\bmod12\}$. This subgroup has index $12$, with coset representatives $h(0,0,k)$, $0\le k\le11$, and it is generated by $c=[T_1,T_2^2]=h(1,0,-6)$ and $c'=T_2cT_2^{-1}=h(0,1,6)$.
\item $G\cap SL_2(\ZZ)=SL_2(\ZZ)'$, the commutator subgroup, which is free of rank $2$ and of index $12$.
\item The abelianisation of $G_{\ZZ}$ is $(\ZZ/12)^2$, via $g\mapsto(\operatorname{ab}(M),\varphi(h))$, and $[G_{\ZZ},G_{\ZZ}]=SL_2(\ZZ)'\ltimes H'$.
\item $G$ contains the principal congruence subgroup of level $12$ and no principal congruence subgroup of level $4$ or $6$. For every $D\ge1$ the image of $G$ in $GL_4(\ZZ/D)$ has index $\gcd(D,12)$ in the image of $G_{\ZZ}$.
\item $G$ is Zariski dense in $P=SL_2\ltimes\mathrm{Heis}$, the algebraic stabiliser of $\gamma$, and it is an arithmetic subgroup of $P$, since it has finite index in $P(\ZZ)=G_{\ZZ}$.
\end{enumerate}
\end{theorem}
\status{Proved here; new. Checked by computer (\note{s6\_monodromy\_checks.py}, $24$ checks; \note{ver52\_referee\_claims.py}, $22$ checks, including that every one of the $165{,}888$ elements of $G$ modulo $12$ has $\Phi=0$), and confirmed by the referee's constructive membership algorithm.}
\begin{proof}
\emph{(a), $\Phi$ is a homomorphism.} $\varphi$ is a homomorphism on $\mathrm{Heis}(\ZZ)$ because $q$ is a quadratic refinement of the mod-$2$ symplectic form $\omega(x,y)=p_xr_y-p_yr_x$: $q(x+y)\equiv q(x)+q(y)+\omega(x,y)\pmod2$, and therefore $6q(x+y)\equiv6q(x)+6q(y)+6\omega(x,y)\pmod{12}$. Moreover $\varphi$ is invariant under the Levi action, since $q$ equals $1$ exactly on the three nonzero vectors of $\Ff_2^2$ and so $q(Mx)\equiv q(x)$. Writing $L_{M_1}h_1\cdot L_{M_2}h_2=L_{M_1M_2}\cdot(L_{M_2}^{-1}h_1L_{M_2})\cdot h_2$ gives $\Phi(g_1g_2)=\Phi(g_1)+\Phi(g_2)$.

\emph{(a), $G\subseteq\ker\Phi$ and a lower bound.} By (b), $\Phi(T_1)=4+(2-6)=0$ and $\Phi(T_2)=9+(-3-6)=0$, so $G\subseteq\ker\Phi$. Since $\Phi(t_\gamma)=1$, $\ker\Phi$ has index $12$; hence $[G_{\ZZ}:G]\ge12$ and $G\cap\mathrm{Heis}(\ZZ)\subseteq\ker\varphi=H'$.

\emph{(a), $G=\ker\Phi$.} The words $c^pc'^r[c,c']^m$ equal $h(p,r,-6p+6r+6pr+12m)$, and $-6p+6r+6pr\equiv6q(p,r)\pmod{12}$, so these words exhaust $H'$ and $H'=\langle c,c'\rangle\subseteq G$. Since $R$ and $S$ generate $SL_2(\ZZ)$ ($S(SR)S^{-1}=T$), $G$ maps onto $SL_2(\ZZ)$. Therefore $[G_{\ZZ}:G]=[\mathrm{Heis}(\ZZ):G\cap\mathrm{Heis}(\ZZ)]\le[\mathrm{Heis}(\ZZ):H']=12$, and $G=\ker\Phi$.

(b) is a direct computation, and (c) follows from (a) and the description of the centre.

(d) $H'$ is $\ker\varphi$, and $h(0,0,k)h(p,r,s)=h(p,r,s+k)$ gives the coset representatives.

(e) $L_M\in G$ exactly when $\operatorname{ab}(M)=0$. The commutator subgroup $SL_2(\ZZ)'$ has index $12$ and is torsion-free: every element of finite order in $SL_2(\ZZ)$ is conjugate to a power of $S$ or of $-R$, and these powers have the $\operatorname{ab}$-values $9,6,3$ and $10,8,6,4,2$, all nonzero. A torsion-free subgroup of the amalgam $\ZZ/4*_{\ZZ/2}\ZZ/6$ acts freely on its Bass--Serre tree and is free \cite[Ch.~I]{SerreTrees}; its rank follows from the Euler characteristic, $12\cdot(-\frac1{12})=-1=1-\mathrm{rank}$.

(f) The map is a homomorphism by the computation in (a), applied to each coordinate, and it is onto, since $L_T\mapsto(1,0)$ and $t_\gamma\mapsto(0,1)$. Its kernel $SL_2(\ZZ)'\ltimes H'$ is generated by commutators: those of Levi elements give $SL_2(\ZZ)'$; the commutators $[L_M,h(x,s)]=h((M-I)x,-6\omega(Mx,x))$ have projections $(M-I)x$ that span $\ZZ^2$, for example $(T-I)(0,1)=(1,0)$ and $(T^{\mathsf T}-I)(1,0)=(0,1)$; and the commutators of these give $h(0,0,12)$. By Proposition~\ref{prop:heis} below the subgroup they generate is $H'$.

(g) If $g\equiv I\pmod{12}$, then $\operatorname{ab}(M)=0$, because $\operatorname{ab}$ factors through $SL_2(\ZZ/12)$ (checked by a breadth-first search over all $1152$ elements), and $\varphi(h)=0$. On the other hand $h(0,0,4)$ and $h(0,0,6)$ lie in the principal congruence subgroups of levels $4$ and $6$ and have $\Phi=4$ and $6$. For the index, let $K_D$ be the principal congruence subgroup of level $D$; then $[G_{\ZZ}\bmod D:G\bmod D]=[G_{\ZZ}:GK_D]=12/|\Phi(K_D)|$. Since $t_\gamma^D\in K_D$, $\Phi(K_D)$ contains $\gcd(D,12)\ZZ/12$. Conversely, $\Phi$ modulo $d=\gcd(D,12)$ factors through reduction modulo $d$: $SL_2(\ZZ/12)=SL_2(\ZZ/4)\times SL_2(\ZZ/3)$ with $SL_2(\ZZ/3)^{\mathrm{ab}}=\ZZ/3$, and $6q(p,r)\bmod d$ depends only on $p,r\bmod2$.

(h) $h(2,0,0)$, $h(0,2,0)$ and $h(0,0,12)$ lie in $H'\subseteq G$. Their cyclic groups are Zariski dense in the three one-parameter subgroups that generate $\mathrm{Heis}$, so the Zariski closure of $G$ contains $\mathrm{Heis}$. Its Levi image is closed and contains $SL_2(\ZZ)$, which is dense in $SL_2$.
\end{proof}

In particular $G$ is not a thin group in the sense of Sarnak \cite{SarnakThin}: it has finite index in the integral points of its Zariski closure. This places the $S^6$ period system among the arithmetic cases of rank-four symplectic monodromy; Brav and Thomas exhibit hypergeometric monodromy groups in $Sp(4,\ZZ)$ of the other, thin kind \cite{BravThomas}.

\subsection{The Heisenberg part, and the twist}

\begin{proposition}\label{prop:heis}
\begin{enumerate}[label=(\alph*)]
\item The quadratic refinements of $\omega$ modulo $2$ on $\Ff_2^2$ are the four functions $\alpha p^2+pr+\beta r^2$ ($\alpha,\beta\in\Ff_2$). Three have Arf invariant $0$, and $q=p^2+pr+r^2$ has Arf invariant $1$. $SL_2(\Ff_2)$ permutes the three even refinements and fixes $q$.
\item Let $N\subseteq\mathrm{Heis}(\ZZ)$ be a subgroup that is normalised by the Levi group and maps onto $\ZZ^2$ under $(p,r)$. Then $N=N_k:=\{h(p,r,s):s\equiv6q(p,r)\bmod k\}$ for a unique divisor $k$ of $12$. Conversely, every $N_k$ is such a subgroup, of index $k$.
\item $H'=N_{12}$ is the smallest of them, and $H'=[G_{\ZZ},\mathrm{Heis}(\ZZ)]=[G_{\ZZ},G_{\ZZ}]\cap\mathrm{Heis}(\ZZ)$.
\item On $H(\ZZ)=\{h(p,r,6\kappa)\}$, $H'$ is the kernel of the character $v_H(h(p,r,6\kappa))=(-1)^{p+r+pr+\kappa}$.
\end{enumerate}
\end{proposition}
\status{Proved here; new.}
\begin{proof}
(a) is a finite check. (b) Lifts of $(1,0)$ and $(0,1)$ in $N$ have commutator $h(0,0,12)$, so $N$ meets the centre in $k\ZZ$ for some $k\mid12$, and the fibre of $N$ over $x\in\ZZ^2$ is $\sigma(x)+k\ZZ$. The function $\lambda(x)=\sigma(x)-6q(x)\bmod k$ is additive, by the group law and (a), and it is Levi-invariant; so it vanishes on every $(M-I)\ZZ^2$, and these span $\ZZ^2$. Hence $\lambda=0$. (c) The subgroup $[G_{\ZZ},\mathrm{Heis}(\ZZ)]$ is Levi-normal, maps onto $\ZZ^2$, contains $h(0,0,12)$ and lies in $H'$, so it is $N_{12}=H'$. (d) With $s=6\kappa$, the condition $s\equiv6q(p,r)\pmod{12}$ says $\kappa\equiv q(p,r)\equiv p+pr+r\pmod2$.
\end{proof}

Two readings of Proposition~\ref{prop:heis} connect the Heisenberg part to established theory. Quadratic refinements of the mod-$2$ intersection form of a surface are its spin structures, with the Arf invariant as their parity \cite{Johnson}; on a Riemann surface they are its theta characteristics \cite{AtiyahSpin}. On a torus the odd spin structure is the one fixed by the mapping class group $SL_2(\ZZ)$, and (b) shows that $SL_2$-invariance forces it. The character $v_H$ of (d) is the Heisenberg multiplier of the odd Jacobi theta function $\vartheta$: Gritsenko and Hulek show that $\vartheta\in J_{1/2,1/2}(v_\eta^3\times v_H)$ \cite[(1.4)--(1.5)]{GritsenkoHulek}. So the Heisenberg part of the monodromy group is the kernel of the odd theta multiplier.

\begin{proposition}\label{prop:twist}
For $k\in\ZZ/12$ put $G_k=\{L_Mh:\varphi(h)+k\operatorname{ab}(M)\equiv0\}$. The $G_k$ are twelve distinct normal subgroups of index $12$. Each has Levi image $SL_2(\ZZ)$ and Heisenberg part $H'$, and $G=G_1$. The extension $1\to H'\to G\to SL_2(\ZZ)\to1$ does not split, whereas $G_0=SL_2(\ZZ)\ltimes H'$ does.
\end{proposition}
\status{Proved here; new.}
\begin{proof}
Distinctness: $L_Th(0,0,s)\in G_k$ exactly when $s\equiv-k$. Non-splitting: a splitting $\sigma$ would give an involution $\sigma(-I)\in G$ over $-I$ that commutes with an element $\sigma(R)\in G$ over $R$. Since $(L_{-I}h(x,s))^2=h(0,2s)$, the involutions of $G$ over $-I$ are the $L_{-I}h(x,0)$ with $\Phi=6-6q(x)\equiv0$, that is, with $x\not\equiv0\pmod2$. Conjugation by $L_Rh(y,t)$ sends $L_{-I}h(x,0)$ to $L_{-I}h(R(x-2y),0)$, so commuting would require $Rx\equiv x\pmod2$; but $R$ has no nonzero fixed vector on $\Ff_2^2$.
\end{proof}

So the pair \emph{index $12$ in the centre, odd refinement modulo $2$} determines the Heisenberg part $H'$, and the group $G$ itself is the twist of $H'$ by the abelianisation of the Levi part.

\subsection{The marked peripheral character of the higher-torus update}

The higher-torus update of the record (its file \fn{15\_}, §3.2 and Theorem~6.1; see also \fn{14\_}, §2) works with the peripheral matrices
\[
C_1=\begin{pmatrix}-1&-1\\1&0\end{pmatrix},\qquad C_2'=\begin{pmatrix}0&1\\-1&0\end{pmatrix},\qquad C_0=\begin{pmatrix}1&1\\0&1\end{pmatrix}.
\]
It defines the marked peripheral character $\chi=\nu\circ\operatorname{ab}:SL_2(\ZZ)\twoheadrightarrow\ZZ/12$ by $\chi(C_0)=-1$, and shows that $(\chi(C_0),\chi(C_1),\chi(C_2'))=(-1,4,-3)$.

\begin{theorem}[the peripheral character]\label{thm:periph}
$C_1=\pi(T_1)^{\mathsf T}$, $C_2'=\pi(T_2)^{\mathsf T}$ and $C_0=T$, where $\pi$ denotes the Levi part, and $\chi=-\operatorname{ab}$. The triple $(-1,4,-3)$ equals $(\operatorname{ab}\pi(T_0),\operatorname{ab}\pi(T_1),\operatorname{ab}\pi(T_2))$. For every $g=L_Mh(p,r,s)$ in the monodromy group $G$,
\[
\chi(M)\equiv s-6(p^2+pr+r^2)\pmod{12}.
\]
\end{theorem}
\status{Proved here; new.}
\begin{proof}
The matrix identities are direct. $\chi$ and $-\operatorname{ab}$ are homomorphisms that agree on $T$, which generates the abelianisation, so they are equal. The triple follows from Theorem~\ref{thm:mono}(b), and the last formula is $\Phi(g)=0$ rewritten.
\end{proof}

So the character that the update attaches to the peripheral monodromy is exactly the Levi half of the character whose kernel is the monodromy group. On $G$ it is determined by the Heisenberg coordinates, and the monodromy group is the locus where the two halves cancel. The same update exhibits a Lie-framed torus whose spin quadratic form is $q(x,y)=x+xy+y$ over $\Ff_2$, with Arf invariant one (\fn{15\_}, (147); \fn{14\_}, §1); since $x+xy+y\equiv x^2+xy+y^2\pmod2$, this is the odd refinement of Proposition~\ref{prop:heis}(a). Whether that torus's first homology is the lattice $\langle\bar u,\bar w\rangle$ of Proposition~\ref{prop:heis} is Question~\ref{q:arfone}.

\subsection{The orbits of the \texorpdfstring{$(3,4,\infty)$}{(3,4,infinity)} covers at every level}

The Erdős--Straus workbench (its package ES-09) uses the dual matrices $A_j=(T_j^{-1})^{\mathsf T}$, acting on $\Lambda=\operatorname{Hom}(V,\ZZ)$ with dual basis $(\hat\gamma,\hat u,\hat w,\hat\delta)$, to build covers from the orbits on the affine hyperplane
\[
H_D=\{\lambda\in\Lambda/D\Lambda:\ \text{the $\hat\gamma$-coordinate of $\lambda$ is }1\}.
\]
Write $\lambda=\hat\gamma+b\hat u+c\hat w+d\hat\delta$ and $e=\gcd(D,6)$.

\begin{theorem}[the orbits at every level]\label{thm:orbitsS6}
For every $D\ge1$ the orbits of $G$ on $H_D$ coincide with the orbits of $G_{\ZZ}$. The orbit of $\lambda$ is $\{\lambda'\in H_D:\gcd(b',c',e)=\gcd(b,c,e)\}$, of size $D^3\cdot\#\{v\in(\ZZ/e)^2:\gcd(v,e)=\gcd(b,c,e)\}/e^2$.
\end{theorem}
\status{Proved here; new. Checked as full orbit partitions for all $D\le14$ (\note{ver52\_referee\_claims.py}) and by the referee for all $D\le24$.}
\begin{proof}
\emph{The dual action.} The dual action of $h(p,r,s)$ is $(b,c,d)\mapsto(b+6r,\,c-6p,\,d-s-pb-rc)$, and $L_M$ acts on $(b,c)$ by $M^{-\mathsf T}$ and fixes $d$.

\emph{The orbits of $G_{\ZZ}$.} The coordinate $d$ can be moved freely, and $(b,c)$ matters only modulo $6(\ZZ/D)^2$. So the $G_{\ZZ}$-orbits correspond to the $SL_2(\ZZ/e)$-orbits on $(\ZZ/e)^2$, which are classified by content, since $SL_2(\ZZ)$ maps onto $SL_2(\ZZ/e)$.

\emph{$G$ has the same orbits.} $G$ is normal and $G_{\ZZ}=G\langle t_\gamma\rangle$ with $t_\gamma$ central, so it suffices to show $t_\gamma\lambda\in G\lambda$. At a point with $(b,c)=(g,0)$ the Levi element $L_{(T^{\mathsf T})^{-1}}$ fixes $\lambda$ and has $\Phi=1$, because $(T^{\mathsf T})^{-1}=STS^{-1}$; hence $t_\gamma L_{(T^{\mathsf T})^{-1}}^{-1}\in G$ moves $\lambda$ exactly as $t_\gamma$ does. Every point is $G_{\ZZ}$-equivalent to one with $c=0$, and normality and centrality transport the conclusion to every point.
\end{proof}

\begin{corollary}\label{cor:orbitsS6}
For $D=2^k$, $k\ge1$, there are exactly two orbits, of sizes $D^3/4$ and $3D^3/4$. For odd $D$ there is one orbit if $3\nmid D$, and there are two, of sizes $D^3/9$ and $8D^3/9$, if $3\mid D$. For $6\mid D$ there are four orbits; at $D=6$ their sizes are $6$, $18$, $48$ and $144$, and at $D=12$ they are $48$, $144$, $384$ and $1152$. In general the orbits at $D=2^km$, $m$ odd, are the products of those at $2^k$ and at $m$.
\end{corollary}
\status{Proved here; new.}
\begin{proof}
Count the vectors of $(\ZZ/e)^2$ by content, as in the Erdős--Straus paper \cite[Corollary~5.2]{ESPaper}, and multiply by $D^3/e^2$; the product statement is the Chinese remainder theorem.
\end{proof}

This settles the question, raised when these covers were examined at even levels $D=2^k$ with $k\le5$, whether the two-orbit structure persists for all $k$ and whether mixed levels decompose as products.

\subsection{What the period system does not have}\label{ssec:mononeg}

Three statements of this section and of the manuscript are negative in form, and each is proved.
\begin{itemize}
\item \emph{The extension does not split} (Proposition~\ref{prop:twist}). No subgroup of $G$ maps isomorphically onto $SL_2(\ZZ)$: an involution over $-I$ would have to commute with an element over $R$, and $R$ has no nonzero fixed vector on $\Ff_2^2$.
\item \emph{The level is exactly $12$} (Theorem~\ref{thm:mono}(g)). $G$ contains no principal congruence subgroup of level $4$ or $6$, since $h(0,0,4)$ and $h(0,0,6)$ lie in them and have $\Phi=4$ and $6$.
\item \emph{No polarised limit mixed Hodge structure at the cusp} (the manuscript's Remark~3.23). With $N$ as above, the form $b_N(x,y)=Q_0(x,Ny)$ on $V/\ker N=\langle w,\delta\rangle$ has $b_N(w,w)=Q_0(w,-u)=6$, $b_N(\delta,\delta)=Q_0(\delta,\gamma)=-1$ and $b_N(w,\delta)=b_N(\delta,w)=0$. So $b_N=\operatorname{diag}(6,-1)$ is indefinite, and $(V,Q_0,N)$ does not support a polarised limit mixed Hodge structure.
\end{itemize}
These results say which standard tools apply as they stand: methods that need a split extension or a polarised limit do not, while methods for congruence subgroups of level $12$ and for arithmetic groups do.
